\chapter{k-Nearest Neighbours}
\labinfo{01\_KNN}{knn.py}{UCI Red Wine Quality: 1,599 samples, 11 chemical features; quality grouped into three bands}{Classification}{Test accuracy 68.1\%; macro F1 0.660}

\section{Method}
The response was grouped as low (quality at most 5), medium (6), and high (at
least 7). After a stratified 70/30 split, five-fold cross-validation selected
$k$ and the distance metric. Standardized Manhattan distance with $k=1$ had the
best mean CV accuracy (65.3\%).

\section{Results}
On the held-out 479 wines, accuracy was 68.1\% and macro F1 was 0.660. Scaling
improved accuracy over the unscaled comparison (62.0\%). The per-class F1 scores
were 0.740, 0.642, and 0.597, respectively; the high-quality category remained
the hardest class to identify.

\begin{figure}[H]\centering
\includegraphics[width=.78\linewidth]{\figpath{01_KNN}{cv_accuracy_vs_k.pdf}}
\caption{Cross-validation accuracy for candidate neighbourhood sizes.}
\end{figure}

\section{Python source}
\lstinputlisting[caption={Complete k-NN implementation},label={lst:lab01}]{01_KNN/knn.py}
